\scanpage{8}

\mspage{4}

Bures le 12.8.1964

Mon cher Serre,

Je commence à avoir de faibles lumières sur la façon de faire \inserted{d}es VA avec les cycles al\inserted{g}ébriquement équivalents à zéro d'une variété projective non singulière $X$ (sur un corps alg clos $k$). Si $X$ est connexe de dimension $n$, je sais associer à chaque entier $i$ \inserted{c}ompris entre $\unsure{0|1}$ et $n$ une VA $J^i(X)$, jouant le rôle d'une ``jacobienne intermédiaire'' au sens de Weil, correspondant à des classes de cycles de codimension $i$ (pour une relation d'équivalence que je ne sais pas bien identifier pour l'instant, mais qui mérite certainement d'être appelée ``équivalence d'Albanese''). Je suis moralement sûr que ce sont ``les bonnes'', bien que je n'ai pas prouvé grand-chose dessus pour l'instant; par contre, j'ai des conjectures en pagaille. Je te signale seulement que $J^1=\operatorname{Pic}^{0}$ et $J^n=\operatorname{Alb}^{0}$, que $J^i$ et $J^{n+1-i}$ sont canoniquement duales l'une de l'autre, que $\dim J^i\leq\frac12 b_{2i-1}$ (nombre de Betti) -- de façon plus précise $T_{\ell}(J^i)$ est un quotient d'un sous-module de $H^{2i-1}(X,\uline{Z}_{\ell}(i))$ (et quand on sera plus savant, on saura prouver que c'est même un sous-module dudit), en caractéristique $0$ on peut également remplacer la cohomologie de Weil par la cohomologie de Hodge $H^i(X,\uline{\Omega}^{i-1}_X)$ ...

\margin{conjecturalement $\uline{t}_{J^i}\unsure{\subset|\in}H^{\unsure{i}}(X,\uline{\Omega}^{i-1}_X)$}
\margin{\deleted{$H^{\unsure{1|i}}(J^i)\simeq$ \unsure{Alg. de Hodge}}}

Moyennant un certain nombre de choses non prouvées (qui risquent d'ailleurs d'être vaches) le théorème de positivité de Hodge se ramène à l'énoncé suivant, que je ne sais pas prouver, et que j'ai envie de te communiquer car il est bien terre à terre. Moralement, il s'agit de prouver que dans le cas où $\dim X=2m-1$, l'autodualité de $J^m$ (qui s'exprime par une classe de correspondance divisorielle sur $J^m\mathbin{x}J^m$ \deleted{anti} \uline{symétrique}, donc provenant -- du moins modulo le facteur 2--

\margin{\unsure{Ça veut dire} qu'il faut tenir compte de la décomposition primitive....}

\mspage{7}
